\section*{ACTIVIDADES A DESARROLLAR DESPUÉS DE UN ACCIDENTE}

\subsection*{OBJETIVO}
Poner en práctica la guía de las acciones que hay que tomar en caso de que ocurra un accidente que afecte la Seguridad Física y Radiológica durante el transporte y operación de los Equipos Radiográficos, analizando las causas que puedan dar lugar a que ocurra éste.

\subsection*{MATERIAL}
\begin{itemize}
  \item Detector de Radiaciones tipo Geiger-Müller.
  \item Equipo Portátil Detector de Radiación es con Alarma Sonora.
  \item Contenedor de Gammagrafía (Con todos sus Accesorios).
  \item Fuente simulada.
  \item Juego de Material de Acondicionamiento y Señalización.
\end{itemize}

\subsection*{DESARROLLO}
Se supone que ha ocurrido un accidente automovilístico en donde el técnico y el auxiliar no han sufrido daño físico de consecuencia y que, a causa de éste, han dejado momentáneamente sin vigilancia el equipo radiográfico y el vehículo en que lo transportan, el tiempo suficiente como para que una persona ajena a la empresa se lleve el contenedor con fuente radiactiva (robo).

En cuanto al técnico o el auxiliar se dan cuenta de la falta del contenedor, se procederá de la siguiente manera:
\begin{enumerate}
  \item Se dará parte de los hechos ocurridos a las oficinas centrales y/o al Encargado de Seguridad Radiológica a la brevedad posible, para poder recibir instrucciones al respecto.
  \item Se efectuará una indagación alrededor del lugar donde se supone fue robada o perdida la fuente, la cual permita recuperarla.
\end{enumerate}